\section{Frey 曲线的构造与性质}
\label{sec:ch10-frey}
% ============================================================================

\providecommand{\rhobar}{\bar{\rho}}

1985 年，Frey 提出的想法之所以震撼数论界，
并不是因为他已经证明了什么，
而是因为他指出：
如果费马方程真有解，
那么这个“初等整数方程”会自动制造出一条性质近乎病态的椭圆曲线。
它的坏约化极其稀薄，
判别式的奇素数部分却又几乎是一个完整的 $p$ 次幂；
换句话说，
它在导子里只留下最少量的信息，
却在判别式里留下远超正常规模的幂次。
这种反差正是后来水平下降得以发生的根源。

\textbf{我们遇到了什么问题？}
假设
\[
  a^p+b^p=c^p
\]
有一个非平凡原始解，
其中 $p\ge 5$ 为素数，$a,b,c\in\Z$ 两两互素，且 $abc\neq 0$。
单看这个等式本身，
我们很难从中直接榨出全局矛盾。
它并不会自然长成同余式、类群计算或单位方程。
Frey 的思路是：
不要在原方程里打转，
而是把它送到椭圆曲线世界中，
让判别式、导子与 Galois 表示替我们说话。

\textbf{核心思想是什么？}
把假想的费马解塞进一个 Weierstrass 方程：
\[
  E_{a,b,c}\colon y^2=x(x-a^p)(x+b^p).
\]
这条曲线的三个根 $0,a^p,-b^p$ 彼此距离非常特殊，
因为它们的差分别是 $a^p,b^p,c^p$。
于是判别式里出现的是
\[
  (a^pb^pc^p)^2,
\]
即各奇素数以 $2p$ 次幂进入。
导子却只记住哪些素数真的坏，
而不记住它们出现了多少次。
因此，Frey 曲线的判别式与导子之间形成一种极不寻常的“瘦导子、大判别式”组合。
Ribet 定理正是利用这一点，把级别一步步压下去。

\textbf{引入之后能得到什么？}
一旦 Frey 曲线被构造出来，
我们就可以把费马方程翻译成三条可操作的算术信息：
第一，它是一条半稳定椭圆曲线，因而落入 Wiles 定理的射程；
第二，对每个奇素数 $q\mid abc$，曲线在 $q$ 处是乘法约化，且最小判别式的指数被 $p$ 整除，
所以 mod $p$ 表示在这些地方看不见惯性；
第三，这正好满足 Ribet 水平下降的局部条件。
Frey 曲线不是证明的终点，
而是把一个 Diophantine 假设送进现代模性机器的入口。

\begin{figure}[ht]
\centering
\begin{tikzpicture}[>=Stealth, font=\small, node distance=1.25cm]
  \node[draw, rounded corners, fill=red!8, minimum width=5.2cm, minimum height=0.95cm] (flt) {假设 $a^p+b^p=c^p$ 有原始解，$p\ge 5$};
  \node[draw, rounded corners, fill=orange!14, below=of flt, minimum width=5.4cm, minimum height=0.95cm] (frey) {构造 $E_{a,b,c}: y^2=x(x-a^p)(x+b^p)$};
  \node[draw, rounded corners, fill=yellow!16, below=of frey, minimum width=5.8cm, minimum height=0.95cm] (disc) {判别式奇部分是 $p$ 次幂，导子只记住平方自由部分};
  \node[draw, rounded corners, fill=green!10, below=of disc, minimum width=5.4cm, minimum height=0.95cm] (rep) {得到 $\rhobar_{E,p}$，且在每个 $q\mid abc$ 处无分歧};
  \node[draw, rounded corners, fill=blue!8, below=of rep, minimum width=5.2cm, minimum height=0.95cm] (lvl) {若曲线模，则可反复水平下降到级别 $2$};
  \node[draw, rounded corners, fill=gray!14, below=of lvl, minimum width=3.8cm, minimum height=0.95cm] (zero) {$\Sk_2(\GammaO(2))=0$};

  \draw[->, very thick] (flt) -- node[right] {Frey} (frey);
  \draw[->, very thick] (frey) -- node[right] {算术性质} (disc);
  \draw[->, very thick] (disc) -- node[right] {Tate 模} (rep);
  \draw[->, very thick] (rep) -- node[right] {Ribet} (lvl);
  \draw[->, very thick] (lvl) -- node[right] {维数公式} (zero);
\end{tikzpicture}
\caption{Frey 构造的核心逻辑：假想费马解先制造一条半稳定椭圆曲线，再制造一份在坏素数处异常“安静”的 residual 表示。}
\label{fig:ch10-frey-construction}
\end{figure}


下面固定标准归一化：
假设 $(a,b,c)=1$，$p\ge 5$ 为奇素数，且交换 $a,b$ 后可令 $2\mid b$。
定义
\[
  E_{a,b,c}\colon y^2=x(x-a^p)(x+b^p).
\]

\begin{proposition}[Frey 曲线的判别式]
\label{prop:ch10-frey-discriminant}
对上式的整模型，有
\[
  \Delta(E_{a,b,c})=16\,a^{2p}b^{2p}c^{2p}=16(abc)^{2p}.
\]
进一步地，在原始解且 $2\mid b$ 的标准最小化之后，最小判别式满足
\[
  \Delta_{\min}(E_{a,b,c})=2^{-8}(abc)^{2p}.
\]
特别地，其奇素数部分是一个完全 $p$ 次幂：
\[
  \prod_{q\text{ odd}} q^{v_q(\Delta_{\min})}
  =\Bigl(\prod_{q\text{ odd}} q^{2v_q(abc)}\Bigr)^p.
\]
\end{proposition}

\begin{proof}
把方程写成
\[
  y^2=(x-e_1)(x-e_2)(x-e_3),
  \qquad e_1=0,\ e_2=a^p,\ e_3=-b^p.
\]
对这种三根形式，判别式公式是
\[
  \Delta=16(e_1-e_2)^2(e_1-e_3)^2(e_2-e_3)^2.
\]
代入可得
\[
  \Delta=16(a^p)^2(b^p)^2(a^p+b^p)^2=16a^{2p}b^{2p}c^{2p}.
\]
这证明了整模型的公式。

当 $(a,b,c)=1$ 且 $2\mid b$ 时，
奇素数处该模型已经最小；
因此对每个奇素数 $q$ 都有
\[
  v_q(\Delta_{\min})=2p\,v_q(abc).
\]
在素数 $2$ 处，需要做一次标准的最小化。
Tate 算法给出：
在上述归一化下，最小模型的判别式比当前模型少去 $2^{12}$，
于是
\[
  \Delta_{\min}=2^{-8}(abc)^{2p}.
\]
关键点不在于 $2$ 处的精确指数，
而在于所有奇素数指数都被 $p$ 整除。
这正是后续水平下降真正使用的信息。
\end{proof}

\textbf{注。}
很多介绍会同时写下这两个判别式公式，
读者初看容易困惑。
其实它们讨论的是两件不同的事：
$16(abc)^{2p}$ 是眼前这个最直接整模型的判别式，
而 $2^{-8}(abc)^{2p}$ 是把它在 $2$ 处最小化后得到的最小判别式。
对费马大定理而言，
真正决定性的只是“奇素数指数都是 $p$ 的倍数”这一点。

\begin{proposition}[半稳定性与导子]
\label{prop:ch10-frey-conductor}
Frey 曲线 $E_{a,b,c}$ 在每个奇素数 $q\mid abc$ 处都是乘法约化，
在每个奇素数 $q\nmid abc$ 处都是好约化。
此外，在标准归一化下它在 $2$ 处也是半稳定的。
因此
\[
  N(E_{a,b,c})=2\,\mathrm{rad}_{\mathrm{odd}}(abc),
\]
其中
\[
  \mathrm{rad}_{\mathrm{odd}}(abc)=\prod_{q\mid abc,\ q\text{ odd}} q.
\]
\end{proposition}

\begin{proof}
若奇素数 $q\nmid abc$，则上面整模型的判别式在 $q$ 处是单位，
故曲线在 $q$ 处好约化。
现在设 $q\mid abc$ 为奇素数。
由于 $(a,b,c)=1$，$q$ 恰好整除三者之一。
例如若 $q\mid a$，则 $q\nmid b,c$，并且
\[
  c_4=16(a^{2p}+a^pb^p+b^{2p})\equiv 16b^{2p}\not\equiv 0\pmod q.
\]
所以 $v_q(\Delta)>0$ 而 $v_q(c_4)=0$。
对奇素数这正是乘法约化的判别准则。
同理可处理 $q\mid b$ 与 $q\mid c$ 的情形。

半稳定性的困难只剩下 $2$ 处。
在原始解且 $2\mid b$ 的归一化下，
Tate 算法显示 Frey 曲线在 $2$ 处没有加性坏约化，
因此也是半稳定的。
半稳定椭圆曲线在每个乘法约化处的导子指数都等于 $1$，
于是
\[
  N(E_{a,b,c})=2\prod_{q\mid abc,\ q\text{ odd}} q.
\]
\end{proof}

由此，Frey 曲线的 odd conductor 恰好只是 $abc$ 的奇平方自由部分。
这已经非常异常；
更异常的是对应的 residual 表示在这些奇坏素数处还会变得无分歧。

\begin{proposition}[Frey 曲线的 residual 表示为何可降级]
\label{prop:ch10-frey-unramified}
设 $q$ 为奇素数且 $q\mid abc$。
则对 residual 表示
\[
  \rhobar_{E_{a,b,c},p}\colon G_\Q\to \GL_2(\F_p),
\]
有 $\rhobar_{E_{a,b,c},p}$ 在 $q$ 处无分歧。
	a
	herefore 若 $E_{a,b,c}$ 模，
则其 mod $p$ 表示满足 Ribet 水平下降在每个 $q\mid abc$ 处所需的局部条件。
\end{proposition}

\begin{proof}
由 \cref{prop:ch10-frey-conductor}，
$E_{a,b,c}$ 在 $q$ 处是乘法约化。
对乘法约化椭圆曲线，Tate 曲线理论说明：
惯性群在 $E[p]$ 上的作用由参数 $q_E$ 控制，
而最小判别式指数正是 $v_q(q_E)$。
由 \cref{prop:ch10-frey-discriminant}，
有
\[
  v_q(\Delta_{\min})=2p\,v_q(abc),
\]
被 $p$ 整除。
于是惯性在 mod $p$ 后变成平凡，
即 $\rhobar_{E_{a,b,c},p}$ 在 $q$ 处无分歧。
这正是“判别式指数是 $p$ 的倍数”会触发水平下降的精确原因。
\end{proof}

\begin{example}[演示性的数值例子]
\label{ex:ch10-frey-345}
取 $a=3,b=4,c=5$，指数取 $1$ 只是为了演示构造，
并不是费马方程的解。
则
\[
  E_{3,4,5}\colon y^2=x(x-3)(x+4)=x^3+x^2-12x.
\]
其判别式是
\[
  \Delta=16\cdot 3^2\cdot 4^2\cdot 5^2=57600=2^8\cdot 3^2\cdot 5^2.
\]
这说明 Frey 构造本身非常具体：
一组三个整数立刻变成一条椭圆曲线。
真正进入 FLT 时，只是把指数 $1$ 换成奇素数 $p$，
于是奇素数部分变成完整的 $p$ 次幂。
\end{example}

Frey 的论文还没有把这条路完全走完，
但它已经把终点的轮廓照亮了。
Serre 在 1987 年进一步提出他的 $\varepsilon$ 猜想：
Frey 曲线的 residual 表示应该有 Serre 权 $2$、Serre 级别 $2$。
换句话说，
如果它来自模形式，
那就只能来自
\[
  \Sk_2(\GammaO(2)).
\]
而这个空间是零。
Ribet 随后证明的水平下降，
恰恰把这一直觉变成了定理的一半。

